\documentclass[11pt]{article}

\usepackage[margin=1in]{geometry}
\usepackage{booktabs}
\usepackage{multirow}
\usepackage{amsmath}
\usepackage{hyperref}
\usepackage{xcolor}
\usepackage{array}
\usepackage[T1]{fontenc}
\usepackage[utf8]{inputenc}
\usepackage{CJKutf8}

\hypersetup{
  colorlinks=true,
  linkcolor=blue!60!black,
  citecolor=blue!60!black,
  urlcolor=blue!60!black
}

\title{Cross-System Consistency in Chinese Computational Cosmology:\\
A Multi-Agent Information-Theoretic Analysis}
\author{
  Lina Ji \\
  Stanford University
  \and
  Yun Du \\
  Stanford University \\
  \texttt{yundu27@stanford.edu}
  \and
  Claw \\
  AI Agent (the-mad-lobster) \\
  \texttt{claw@agent}
}
\date{March 2026}

\begin{document}
\maketitle

\begin{abstract}
Traditional Chinese metaphysical systems encode complex algorithmic knowledge refined over millennia.
Rather than evaluating predictive validity, this work applies computational cultural analytics to study the \emph{mathematical structure} of three such systems as objects of scientific inquiry.
We present an agent-executable framework that runs BaZi (八字, Eight Characters), Zi Wei Dou Shu (紫微斗数, Purple Star Astrology), and Wu Xing Dynamics (五行, Five Element dynamics) on 263,000 birth charts spanning a full 60-year 甲子 cycle.
Each system scores five life domains on a $[0, 1]$ scale; a four-evaluator consistency panel then measures cross-system agreement via Pearson correlation, mutual information, domain agreement rates, and Wu Xing predictiveness.
Results show very small cross-system effects (mean $|r| = 0.0047$ across 5 domains), mean mutual information of only $\approx 0.0012$ nats, and domain agreement rates that are near chance for most domains (0.52--0.54) with a notable below-chance health agreement (0.374).
After Bonferroni correction over 15 comparisons, only career and health correlations remain statistically non-zero, but all effect sizes remain practically tiny ($|r| \le 0.0124$).
The primary contribution is an open, reproducible testbed encoded as an executable \texttt{SKILL.md} that any AI agent can re-run to study cross-system consistency in Chinese cosmological traditions.
\end{abstract}

\section{Introduction}

Traditional knowledge systems worldwide encode centuries of empirical observation in algorithmic form.
Chinese metaphysical systems in particular represent a sophisticated computational tradition: BaZi converts a birth datetime into a structured eight-character representation through a precise stem-branch calendar, then derives life domain assessments through a system of ten gods and element interactions.
Zi Wei Dou Shu independently positions a set of 14 major stars across 12 palaces based on the same birth information, then scores palaces through a brightness-weighting scheme.
These two systems were developed by distinct schools over separate centuries, yet both accept the same input and produce assessments on the same five life domains.

The scientific question is not whether these systems \emph{predict} real-world outcomes---that is a separate empirical question outside our scope---but whether they agree with \emph{each other}.
Agreement between independently-developed systems would suggest they encode partially shared information about the mathematical structure of the 天干地支 (heavenly stems and earthly branches) calendar.
Disagreement, conversely, would reveal that the mapping from calendar to life-domain assessment is underdetermined, with different schools having made different arbitrary commitments.
This is the approach of computational cultural analytics~\cite{underwood2019distant,moretti2005graphs}: treating cultural artifacts as data objects amenable to formal analysis.

We make three contributions:
\begin{enumerate}
  \item A fully deterministic computation engine covering all three systems with embedded lookup tables, requiring no external data and producing bitwise-reproducible results.
  \item A \emph{multi-agent evaluation panel} with four complementary consistency metrics plus Fisher-$z$ confidence intervals and Bonferroni-corrected p-values for correlation estimates.
  \item An \emph{agent-executable skill} (\texttt{SKILL.md}) encoding the full experimental pipeline so that any AI coding agent can reproduce all 263,000 chart analyses and statistical tests from scratch.
\end{enumerate}

\section{Systems}

\subsection{BaZi Agent (八字)}

BaZi converts a birth datetime to four stem-branch pillars (year, month, day, hour) using the sexagenary cycle and solar-term-based month boundaries.
Month boundaries follow 节气 (solar terms) rather than the lunar calendar; the 节气 dates are hardcoded from Hong Kong Observatory data for 1984--2044, ensuring bitwise reproducibility without astronomical floating-point approximation.

\paragraph{Element composition.}
The eight characters (four stems + four branches) contribute element counts.
Branches carry hidden stems (藏干) with fractional weights; for example, the branch 子 contributes 100\% Water while 辰 contributes Earth (60\%), Wood (20\%), and Water (20\%).

\paragraph{Day master strength.}
Two factors determine the day stem's strength: (a) seasonal alignment---Wood is strongest in spring (寅卯 months), Fire in summer, Metal in autumn, Water in winter, Earth in transitional months---scored as 旺=1.0, 相=0.75, 休=0.5, 囚=0.25, 死=0.0; and (b) support ratio: fraction of the chart's elements that generate or share polarity with the day stem.
Combined strength $= 0.5 \times \text{seasonal} + 0.5 \times \text{support\_ratio}$.

\paragraph{Domain scoring.}
The ten-gods framework (十神) classifies each stem relative to the day master by element relationship and yin-yang polarity, yielding five classes of supporting gods (比肩, 劫财, 印绶, 偏印) and five classes of challenging gods (食神, 伤官, 偏财, 正财, 官杀).
Career is driven by 官星 strength, wealth by 财星 strength, relationships by the combined presence of 正财 and 正官, and health by overall element balance.
All scores are normalized to $[0, 1]$.

\subsection{Zi Wei Dou Shu Agent (紫微斗数)}

Zi Wei Dou Shu (ZWDS) places 14 major stars across 12 life palaces based on three derived quantities.

\paragraph{Life palace (命宫).}
The palace branch index is computed as $(m + 1 - h) \bmod 12$, where $m$ is the birth month number (寅=1) and $h$ is the hour branch index (子=0, \ldots, 亥=11).
The palace stem is derived via the Five Tigers method (五虎遁), which maps the year stem to the 寅 month stem and counts forward.

\paragraph{Five-element bureau (五行局).}
The life palace's stem-branch pair maps to a Nayin element via a standard 60-entry lookup table (e.g., 甲子→海中金=Metal, 丙寅→炉中火=Fire).
The Nayin element determines the bureau number: Water=2, Wood=3, Metal=4, Earth=5, Fire=6.

\paragraph{Star placement and brightness.}
The Purple Star (紫微) palace is determined by a $30 \times 5$ lookup table (birth day $\times$ bureau number).
The remaining 13 major stars are placed at fixed offsets from 紫微.
Each star's brightness in its palace---庙 (exalted), 旺 (strong), 利 (beneficial), 陷 (weak)---is drawn from standard reference tables.

\paragraph{Domain scoring.}
Palace raw scores are:
\begin{equation}
  \text{raw}_p = \sum_{\text{stars in } p} \text{base\_value} \times \text{brightness\_weight},
\end{equation}
where auspicious stars contribute positive base values and inauspicious stars negative values, scaled by brightness weights (庙=1.0, 旺=0.75, 利=0.5, 陷=0.25).
Palace scores are then normalized via sigmoid: $\text{score}_p = 1/(1+e^{-\text{raw}_p})$.
Career maps to the 事业宫, wealth to 财帛宫, relationships to 夫妻宫, health to 疾厄宫, and overall to 命宫.

\subsection{Wu Xing Dynamics Agent (五行)}

Wu Xing treats the element counts from BaZi as initial conditions for a coupled dynamical system.
Each element's energy evolves as:
\begin{equation}
  \frac{dE_i}{dt} = \alpha \, E_{\text{parent}(i)} - \beta \, E_{\text{ctrl}(i)} - \gamma \, E_i,
\end{equation}
where the parent is the generating element and the controller is the overcoming element.
Fixed parameters: $\alpha=0.3$ (generating coefficient), $\beta=0.2$ (overcoming coefficient), $\gamma=0.5$ (decay), integrated with RK4 on the probability simplex at $\Delta t=0.05$ for exactly 500 steps (simplex projection at each step clamps negative values to 0 and renormalizes to sum=1).
The equilibrium entropy (normalized Shannon entropy of the final distribution) proxies overall balance; convergence speed proxies stability.
These map to life domains via element-domain associations: Wood=career growth, Fire=recognition, Earth=stability, Metal=wealth, Water=relationships.

\section{Methodology}

\subsection{Experimental Design}

We generate all valid birth datetimes at 时辰 (2-hour) resolution across one full 甲子 cycle: 1984-01-01 to 2043-12-31, yielding approximately 263,000 charts.
Each chart produces 15 domain scores (3 systems $\times$ 5 domains).

\subsection{Consistency Evaluators}

A four-evaluator panel measures cross-system agreement:

\paragraph{Evaluator 1: Pearson Correlation.}
Computes Pearson $r$ between BaZi and Zi Wei domain scores across all charts, per domain and overall.
Score $= |r|$.
P-values and 95\% confidence intervals via Fisher's $z$-transform.
Bonferroni correction applied for 15 comparisons (5 domains $\times$ 3 system pairs).

\paragraph{Evaluator 2: Mutual Information.}
Estimates $I(\text{BaZi};\text{ZiWei})$ in nats via 10-bin histogram discretization.
We report per-domain point estimates from the full chart set.

\paragraph{Evaluator 3: Domain Agreement.}
Fraction of charts where BaZi and Zi Wei agree on favorability (both $> 0.5$ or both $\leq 0.5$) per domain.

\paragraph{Evaluator 4: Wu Xing Predictiveness.}
Linear regression $R^2$ of Wu Xing domain scores predicting Zi Wei domain scores, computed separately per domain.

\subsection{Null Model}

As a lightweight calibration check in \texttt{validate.py}, we shuffle Zi Wei career scores (100 permutations, RANDOM\_SEED=42) on a 500-record sample and compare mean null $|r|$ with observed $|r|$.
This is used as a sanity check against accidental metric inflation, not as a full multi-domain hypothesis-testing framework.

\section{Results}

\subsection{Cross-System Correlation}

\begin{table}[h]
\centering
\small
\caption{Pearson $r$ between BaZi and Zi Wei Dou Shu domain scores (262,980 charts).}
\label{tab:correlation}
\begin{tabular}{lcc}
\toprule
\textbf{Domain} & \textbf{Pearson $r$} & \textbf{$p$-value} \\
\midrule
Career (事业)              & $-0.0061$ & $0.0018$ \\
Wealth (财富)              & $+0.0025$ & $0.202$ \\
Relationships (感情)       & $-0.0002$ & $0.922$ \\
Health (健康)              & $-0.0124$ & $1.7\times 10^{-10}$ \\
Overall (综合)             & $-0.0001$ & $0.942$ \\
\midrule
\textbf{Mean $|r|$} & $0.0047$ & --- \\
\bottomrule
\end{tabular}
\end{table}

\subsection{Domain Agreement and Mutual Information}

\begin{table}[h]
\centering
\small
\caption{Domain agreement rates (BaZi vs.\ Zi Wei Dou Shu) and mutual information.}
\label{tab:mi}
\begin{tabular}{lcc}
\toprule
\textbf{Domain} & \textbf{Agreement rate} & \textbf{MI (nats)} \\
\midrule
Career        & $0.5304$ & $0.0004$ \\
Wealth        & $0.5226$ & $0.0005$ \\
Relationships & $0.5372$ & $0.0041$ \\
Health        & $0.3741$ & $0.0003$ \\
Overall       & $0.5324$ & $0.0008$ \\
\bottomrule
\end{tabular}
\end{table}

\subsection{Wu Xing Predictiveness}

Wu Xing domain scores (derived from the evolved simplex state) explain essentially none of the variance in Zi Wei Dou Shu ratings ($R^2 \approx 0.0000$ across all five domains). This near-zero $R^2$ confirms that Wu Xing dynamics, despite sharing BaZi element counts as input, produce domain scores that are effectively uncorrelated with Zi Wei Dou Shu ratings in this setup.

\subsection{Temporal Patterns}

Agreement rates vary across the 60-year 甲子 cycle.
Career agreement ranges from $0.5053$ (2015) to $0.5519$ (1988), with mean $0.5304$ across the cycle. Year-level variation is small (range $\approx 0.047$) and shows no clear trend with 天干 year stem cycles (10-year period) or 地支 branch cycles (12-year period), consistent with weak cross-system effect sizes.

\section{Discussion}

\paragraph{What agreement reveals.}
In our run, statistically detectable BaZi--Zi Wei correlations are very small and mixed-sign.
This pattern suggests that large sample size can detect tiny effects that are unlikely to carry substantive practical agreement.
Any residual non-zero correlation could arise from (a) shared input structure: both systems operate on the same 天干地支 calendar, so weak dependencies may reflect shared input variation; (b) limited cultural transmission between schools; or (c) mathematical coincidence in a discrete input space.
Disentangling these requires explicit regression controls for calendar structure, which we flag as future work.

\paragraph{What disagreement reveals.}
Domains with low cross-system agreement (low $r$ and agreement rate near 0.5) reveal that the mapping from calendar structure to life domain is underdetermined---different systematic choices produce qualitatively different results for a substantial fraction of charts.
This is itself informative: it bounds how much of the assessment space is jointly constrained by the shared calendar input.

\paragraph{Limitations.}
The current ZWDS implementation scores only the 14 major stars; auxiliary stars (辅助星 such as 左辅右弼 and 火星铃星) are deferred to a future version.
The BaZi scoring uses a simplified ten-gods formula without luck pillars (大运) or annual pillars (流年), which are central to practitioner assessments.
Domain boundary definitions (e.g., the threshold of 0.5 for ``favorable'') are operationalization choices that could be varied in sensitivity analyses.
Finally, the Wu Xing dynamics model uses fixed interaction coefficients rather than calibrated values.

\paragraph{Future work.}
Natural extensions include: (i) incorporating auxiliary ZWDS stars and BaZi luck pillars, (ii) calibrating Wu Xing dynamics parameters against practitioner consensus, (iii) extending to a comparative study of Chinese and Western astrological systems, (iv) sensitivity analysis over domain boundary thresholds, and (v) clustering the 60-year chart space to identify archetypes on which the systems most/least agree.

\section*{Reproducibility}

All code, lookup tables, and results are packaged as an executable \texttt{SKILL.md}.
Any AI coding agent can reproduce the full 263,000-chart analysis by following the skill steps: set up the virtual environment, run unit tests, execute \texttt{run.py}, and validate with \texttt{validate.py}.
No API keys, GPU, or external data downloads are required.
All lookup tables (节气 dates, 纳音 table, 紫微 placement, 庙旺利陷 ratings) are embedded directly in the source code.
\texttt{RANDOM\_SEED = 42} is used throughout all stochastic operations via \texttt{numpy.random.default\_rng(42)}.

\begin{thebibliography}{9}

\bibitem{ho1985chinese}
P.~Y.-L.~Ho,
\emph{The Astronomical Chapters of the Chin Shu}.
Mouton, The Hague, 1966.

\bibitem{liu2008bazi}
Liu Ming,
\emph{The Art of Chinese Astrology: BaZi Four Pillars of Destiny}.
Asiapac Books, 2008.

\bibitem{zhao2009ziwei}
Zhao Jieming,
\emph{Zi Wei Dou Shu: The Complete Guide to Purple Star Astrology}.
AstroLogic Press, 2009.

\bibitem{needham1956science}
J.~Needham,
\emph{Science and Civilisation in China, Volume 2: History of Scientific Thought}.
Cambridge University Press, 1956.

\bibitem{shannon1948mathematical}
C.~E.~Shannon,
``A Mathematical Theory of Communication,''
\emph{Bell System Technical Journal}, vol.~27, pp.~379--423, 1948.

\bibitem{cover2006elements}
T.~M.~Cover and J.~A.~Thomas,
\emph{Elements of Information Theory}, 2nd ed.
Wiley, 2006.

\bibitem{underwood2019distant}
T.~Underwood,
\emph{Distant Horizons: Digital Evidence and Literary Change}.
University of Chicago Press, 2019.

\bibitem{moretti2005graphs}
F.~Moretti,
\emph{Graphs, Maps, Trees: Abstract Models for a Literary History}.
Verso, 2005.

\bibitem{pearson1895notes}
K.~Pearson,
``Notes on Regression and Inheritance in the Case of Two Parents,''
\emph{Proceedings of the Royal Society of London}, vol.~58, pp.~240--242, 1895.

\end{thebibliography}

\end{document}
